\documentclass[11pt]{article}

\usepackage[margin=1in]{geometry}
\usepackage{fancyhdr}

\title{\Large PLUS Session Outline \\ \normalsize \textit{MATH 1554:} Linear Algebra}
\author{Ahmed Mohammed}
\date{27th of June, 2023}

% ------- Set up header -------
\renewcommand{\sectionmark}[1]{\markright{#1}} % from https://tex.stackexchange.com/a/317393/227665
\pagestyle{fancy}
\fancyhead[L]{27th of June, 2023}
% \fancyhead[C]{\textit{\nouppercase{\rightmark}}}
\fancyhead[C]{\textit{MATH 1554 PLUS Session Outline}}
\fancyhead[R]{Mohammed \textbf{\thepage}}
% -------------------------------------------------------

\begin{document}
%\maketitle
\thispagestyle{fancy}
\section{\centering Open Activity}
	\begin{enumerate}
		\item \textit{Activity}. The group(s) will be asked a conceptual question and will be given a moment to answer it. Then, the group will be introduced to the 
			main topic of the session.
		\item \textit{Collaboration}. The group(s) will be expected to ponder upon the question together in order to reach to an answer.
		\item \textit{Intended Goal}. The intention of this activity is to serve as a motivation for the topic of this session. It (partially) answers \textit{why} we
			care about it.
	\end{enumerate}

\section{\centering Main Activity}
	\begin{enumerate}
		\item \textit{Activity}. The answer to the previous conceptual problem will be revealed in the form of another question which I'll go through. Students will
			then be asked to answer more questions one by one, with me providing them with enough time for them to answer. I'll then go through each 
			question.
		\item \textit{Collaboration}. During the amount of time the students are pondering upon each question, they are expected to work with each other.
		\item \textit{Intended Goal}. This activity furthermore shows \textit{why} the main topic of the lecture is powerful. It also serves as a practice for 
			answering similar questions found in previous exams.
	\end{enumerate}

\section{\centering Closing Activity}
	\begin{enumerate}
		\item \textit{Activity}. In this activity, students will be lead to other resources that discuss the main topic of this session in more detail, or in an alternative
			fashion. Then, the next session's main topic is ``foreshadowed."
		\item \textit{Collaboration}. \textsc{N/A}
		\item \textit{Intended Goal}. This activity provides students with more resources, and introduces the topic of the next session.
	\end{enumerate}

\newpage
\fancyhead[L]{29th of June, 2023}
\section{\centering Open Activity}
	\begin{enumerate}
		\item \textit{Activity}. The group will be introduced to the main topic.
		\item \textit{Collaboration}. \textsc{N/A}
		\item \textit{Intended Goal}. The intention of this activity is to add a level of intuition into the topic.
	\end{enumerate}

\section{\centering Main Activity}
	\begin{enumerate}
		\item \textit{Activity}. The group of students are asked a single question one-at-a-time and are given time to answer it. Then, I walk through the 
			questions.
		\item \textit{Collaboration}. During the amount of time the students are pondering upon each question, they are expected to work with each other.
		\item \textit{Intended Goal}. This activity is meant to add a level of concreteness into the discussion and to deal with questions they might encounter
			in future exams.
	\end{enumerate}

\section{\centering Closing Activity}
	N/A.

\newpage
\fancyhead[L]{6th of July, 2023}
\section{\centering Open Activity}
	\begin{enumerate}
		\item \textit{Activity}. Students will review the topics covered in the past few lectures
		\item \textit{Collaboration}. \textsc{N/A}
		\item \textit{Intended Goal}. Students will have a similar base-ground understanding of the topics relevant to this session.
	\end{enumerate}

\section{\centering Main Activity}
	\begin{enumerate}
		\item \textit{Activity}. Students will solve a series of problems that pertrain to the topics covered in this session. Namely, the students will be asked
						to solve computational problems when it comes to determining the distance between vectors and subspaces and vectors.
						Students will also be asked to compute the orthogonal projection of a vector onto a subspace, and will be asked conceptual
						questions that'll reinforce there understanding of the first three sections of Chapter 6.
		\item \textit{Collaboration}. Students will be expected to work together in groups. They'll be given an enough amount of time for each problem they
							encounter.
		\item \textit{Intended Goal}. Students will have the opportunity to ask questions and understand their level of comprehension of the topics.
	\end{enumerate}

\section{\centering Closing Activity}
	\begin{enumerate}
		\item \textit{Activity}. We will collectively revise the main topics of the session.
		\item \textit{Collaboration}. \textsc{N/A}
		\item \textit{Intended Goal}. Students will leave the session with an understanding of the topics that they'll have to focus on for the upcoming exam.
	\end{enumerate}

\newpage
\fancyhead[L]{9th of July, 2023 (Double)}
\section{\centering Open Activity}
	\begin{enumerate}
		\item \textit{Activity}. We will summarize the main topics involved in the Least-Square Solution(s) problem.
		\item \textit{Collaboration}.  N/A
		\item \textit{Intended Goal}. This activity will serve as a review for the most recent lesson.
	\end{enumerate}

\section{\centering Main Activity}
	\begin{enumerate}
		\item \textit{Activity}. Students will be asked to work on a worksheet as a review for the upcoming exam. The topics covered in the worksheet include
						lessons from both Chapter 5 and Chapter 6 of the textbook.
		\item \textit{Collaboration}. Students will be expected to work in groups together.
		\item \textit{Intended Goal}. This activity will act as a mean to practice questions for the exam.
	\end{enumerate}

\section{\centering Closing Activity}
	\begin{enumerate}
		\item \textit{Activity}. We will review the main struggle points the students faced in answering the questions.
		\item \textit{Collaboration}. N/A
		\item \textit{Intended Goal}. This activity also acts as a final review for the upcoming exam on Wednesday. It also highlights to the students their 
							potential weakpoints and how they may improve them.
	\end{enumerate}

\end{document}